\subsection{Paths of a prescribed length}

We now show how the graph supplied by
Proposition~\ref{clq:layer-interface} yields a forbidden cycle.
The main step concerns paths whose ends lie in different parts of a
partition.  We first record the elementary longest-path bound needed
when all vertices of the layer have a common tree parent.

\begin{lemma}\label{clq:long-path}
Let $d\ge 2$, and let $F$ be a finite connected graph with
$\delta(F)\ge d$.  Then $F$ contains a path of order at least
$\min\{|V(F)|,2d+1\}$.
\end{lemma}

\begin{proof}
Let $v_1,\ldots,v_r$ be a longest path in $F$.  All neighbors of either
end belong to the path.  Suppose that
$r<\min\{|V(F)|,2d+1\}$.  The two sets
\[
 \{j\in\{1,\ldots,r-1\}:v_jv_r\in E(F)\},\qquad
 \{j\in\{1,\ldots,r-1\}:v_1v_{j+1}\in E(F)\}
\]
have at least $d$ elements each, whereas their common index set has
$r-1\le 2d-1$ elements.  Choose an index $j$ in their intersection.
If $j=1$ or $j=r-1$, the edge $v_1v_r$ closes the path to a cycle.
For $2\le j\le r-2$, the two indicated edges instead give the cycle
\[
 v_1,\ldots,v_j,v_r,v_{r-1},\ldots,v_{j+1},v_1
\]
through all vertices of the path.  Thus in every case there is a
cycle on these $r$ vertices.
Since $r<|V(F)|$ and $F$ is connected, some vertex outside the cycle
has a neighbor on it.  Starting with that vertex and then traversing
the cycle with one edge omitted gives a longer path, a contradiction.
\end{proof}

The next lemma uses only minimum degree, the exclusion of a clique,
and expansion of independent pairs.  Its proof identifies a clique of
possible path ends, shows that this clique attaches to the remaining
graph through one vertex, and uses that attachment to extend a path.

\begin{lemma}\label{clq:coloured-path}
Let $s\ge 4$, and let $H$ be a finite connected nonbipartite graph
containing no $K_s$.  Suppose that
\[
 \delta(H)\ge s,
 \qquad
 |N_H[\{a,b\}]|\ge 2s
 \quad\text{whenever $a,b$ are distinct and nonadjacent.}
\]
For every nonconstant map $\chi:V(H)\to\{0,1\}$ and every integer
$1\le\ell\le 2s-2$, there is a path in $H$ of length exactly $\ell$
whose ends receive different values of $\chi$.
\end{lemma}

\begin{proof}
Fix $\chi$ and $\ell$, and suppose that no such path exists.
There are vertices $y,z$ with $\chi(y)\ne\chi(z)$ and a common
neighbor $x$.  Indeed, if no such pair existed, the color would be
preserved along every walk of length two.  In a connected nonbipartite
graph there is an even walk between any two vertices: if a connecting
path has odd length, first insert an odd closed walk obtained by going
to an odd cycle, traversing it, and returning.  Preservation of the
color along successive pairs of steps would make $\chi$ constant.
This proves the assertion, and $x,y,z$ are distinct because the graph
is simple.

Let $K$ be the component containing $x$ in $H-\{y,z\}$.
Every vertex of $K$ loses at most two neighbors on deleting $y,z$,
and all its other neighbors lie in the same component.  Thus
\begin{equation}\label{clq:component-degree}
 \delta(K)\ge s-2.
\end{equation}
There is no path of order at least $\ell$ starting at $x$ in $K$.
For the first $\ell$ vertices of such a path form a path with some
last vertex $b$.  At least one of $y,z$ has color different from $b$;
prefixing that vertex gives a path of length $\ell$ with differently
colored ends.

Choose a longest path $P$ starting at $x$ in $K$, and write
$W=V(P)$ and $r=|W|$.  All neighbors in $K$ of its last vertex lie
on $P$, so \eqref{clq:component-degree} gives $r\ge s-1$.
If $\ell\le s-1$, this already contradicts the preceding paragraph.
This disposes in particular of $\ell=1,2$.  Henceforth we have
\begin{equation}\label{clq:path-order}
 s\le\ell\le 2s-2,
 \qquad s-1\le r\le\ell-1.
\end{equation}

\medskip
\noindent\emph{The possible last vertices form a clique.}
Define
\[
 E=\{e\in W:\text{there is a path from $x$ to $e$
 whose vertex set is exactly $W$}\}.
\]
Every such path has the maximum possible order among paths starting
at $x$ in $K$.  Its last vertex has no neighbor in $K-W$.
Since vertices of $K$ have no neighbors in
$H-(K\cup\{y,z\})$, we obtain
\begin{equation}\label{clq:endpoint-neighbors}
 N_K(e)\subseteq W,
 \qquad N_H(e)\subseteq W\cup\{y,z\}
 \quad(e\in E).
\end{equation}
If two vertices of $E$ were nonadjacent, their closed neighborhood
in $H$ would have size at most
$r+2\le\ell+1\le 2s-1$, contrary to the hypothesis.
Thus $E$ is a clique.  Also $x\notin E$, because a simple path with
more than one vertex has distinct ends, and $r\ge s-1\ge3$.

We use the standard fixed-end path rotation of P\'osa
\cite[p.~358]{Posa1963}. The endpoint-clique and attachment conclusions
needed here follow from our stated hypotheses. Explicitly, let
$v_1=x,\ldots,v_r=e$ be any path from $x$ with vertex set $W$.
For an edge $ev_j$ with $1\le j\le r-2$, the sequence
\begin{equation}\label{clq:rotation}
 v_1,\ldots,v_j,v_r,v_{r-1},\ldots,v_{j+1}
\end{equation}
is another path from $x$ on exactly $W$.  Its last vertex is
$v_{j+1}$.  The operation is valid also when $j=1$, so the pivot
may be $x$.  For $j=r-1$, the original path itself has last vertex
$v_{j+1}=e$.  Therefore the successor of every neighbor of $e$
along the path belongs to $E$.  Different neighbors have different
successors, and hence
\begin{equation}\label{clq:endpoint-degree}
 d_K(e)\le |E| \qquad(e\in E).
\end{equation}
Combining \eqref{clq:component-degree} with the fact that $E$ is a
clique and $H$ has no $K_s$ gives
\begin{equation}\label{clq:endpoint-size}
 s-2\le |E|\le s-1.
\end{equation}

On every path from $x$ with vertex set $W$, the vertices of $E$ form
a consecutive final segment.  To see this, let $e$ be its last
vertex.  Every earlier vertex of $E$ is adjacent to $e$, since $E$
is a clique.  The rotation then places its successor in $E$ as well.
Thus membership in $E$ is preserved on moving forward along the path.

\medskip
\noindent\emph{The endpoint clique attaches through one vertex.}
Fix one of the paths from $x$ on $W$.  Let $u$ be the predecessor
of its final segment $E$, and let $w_0$ be the first vertex of that
segment.  Both are defined because $E$ is nonempty and $x\notin E$;
in particular, $u\notin E$ and $uw_0$ is an edge.

For $w\in E-\{w_0\}$, reorder the final clique segment so that it
still starts at $w_0$ and now ends at $w$.  Every neighbor of $w$
in $K$ lies on this path by \eqref{clq:endpoint-neighbors}, and
the successor of each neighbor belongs to $E$ by the rotation.
A vertex of the prefix strictly before $u$ has its successor outside
$E$, so it cannot be a neighbor of $w$.  Consequently
\begin{equation}\label{clq:partial-attachment}
 N_K(w)\subseteq (E-\{w\})\cup\{u\}
 \qquad(w\in E-\{w_0\}).
\end{equation}

We claim that $u$ has at least two neighbors in $E$.  There are two
possibilities in \eqref{clq:endpoint-size}.
If $|E|=s-2$, each $w\in E-\{w_0\}$ has at most
$(s-3)+3=s$ possible neighbors in $H$: the other vertices of $E$,
and $u,y,z$.  Minimum degree $s$ forces all these adjacencies.
Since $|E|\ge2$, there is such a $w$, and $u$ is adjacent to both
$w_0$ and $w$.

If $|E|=s-1$, suppose instead that $w_0$ is the only neighbor of
$u$ in $E$.  By \eqref{clq:partial-attachment}, every vertex in
$E-\{w_0\}$ must then see both $y$ and $z$ to have degree at
least $s$.  The vertex $w_0$ is itself a possible last vertex by
the definition of $E$, whether or not it is last in our fixed order.
Thus \eqref{clq:endpoint-degree} gives $d_K(w_0)\le s-1$,
so $w_0$ sees at least one of $y,z$.  That vertex is adjacent to
all of $E$, giving a $K_s$, a contradiction.
We have proved the claim.

Choose a second neighbor $w_1\in E-\{w_0\}$ of $u$.
Keeping the prefix through $u$ fixed, we can enter $E$ at $w_1$
and reorder this clique so that $w_0$ is last.  The same successor
argument used for \eqref{clq:partial-attachment} now applies to
$w_0$ as well.  It follows that
\begin{equation}\label{clq:attachment}
 N_K(E)\setminus E\subseteq\{u\},
 \qquad |N_K(u)\cap E|\ge2.
\end{equation}

There is a vertex $w\in E$ adjacent to both $y,z$.
For $|E|=s-2$, any vertex of $E-\{w_0\}$ already has this
property.  For $|E|=s-1$, suppose none does.  Each vertex of $E$
then has at most one neighbor among $y,z$ and, by
\eqref{clq:attachment}, no neighbor in $K-E$ other than $u$.
To have degree at least $s$, it must be adjacent to $u$.
Thus $E\cup\{u\}$ would be a $K_s$, again a contradiction.

\medskip
\noindent\emph{A path starting in the endpoint clique is long enough.}
We now seek a path of order at least $\ell$ starting at this vertex
$w$ in $K$.  Since both $y$ and $z$ are adjacent to $w$, the same
prefix argument used at $x$ will give the final contradiction.

First, $K$ is not a clique.  Its minimum degree would force a clique
$K$ to have at least $s-1$ vertices.  It cannot have at least $s$
vertices, and if it has exactly $s-1$, minimum degree $s$ in $H$
forces every vertex of $K$ to see both $y,z$.  Then
$K\cup\{y\}$ is a $K_s$.  Hence $K$ has a nonadjacent pair.
Its closed neighborhood in $H$ is contained in $K\cup\{y,z\}$,
so the independent-pair hypothesis gives
\begin{equation}\label{clq:component-size}
 |V(K)|\ge 2s-2.
\end{equation}

The graph $K-E$ is connected.  Otherwise, since $K$ is connected,
each component of $K-E$ would have a neighbor in $E$;
\eqref{clq:attachment} would force every one of these components
to contain $u$.  Moreover,
\[
 |V(K-E)|\ge (2s-2)-(s-1)=s-1\ge3.
\]
Let $R$ be a longest path starting at $u$ in $K-E$, and let $v$
be its last vertex.  Connectedness and the preceding size bound
ensure that $R$ is nontrivial, so $v\ne u$.
By \eqref{clq:attachment}, the vertex $v$ has no neighbor in $E$;
by maximality of $R$, all its neighbors in $K-E$ lie on $R$.
Thus all its neighbors in $K$ lie on $R$, and
\begin{equation}\label{clq:remaining-path-order}
 |V(R)|\ge d_K(v)+1\ge s-1.
\end{equation}

If $|E|=s-2$ and $|V(R)|=s-1$, the vertices $v,w$ are
nonadjacent by \eqref{clq:attachment}, since $v\ne u$.
Their closed neighborhood in $H$ is contained in
\[
 E\cup V(R)\cup\{y,z\}.
\]
These sets are disjoint and have total size
$(s-2)+(s-1)+2=2s-1$, contrary to the hypothesis.
Together with \eqref{clq:endpoint-size} and
\eqref{clq:remaining-path-order}, this proves
\[
 |E|+|V(R)|\ge2s-2\ge\ell.
\]

By \eqref{clq:attachment}, choose a neighbor $e$ of $u$ in $E$
different from $w$.  Traverse all of the clique $E$ in a path
starting at $w$ and ending at $e$, follow the edge $eu$, and then
follow $R$ from $u$, as shown in Figure~\ref{fig:clique-endpoints}.
The vertex sets $E$ and $V(R)$ are disjoint,
so this is a simple path in $K$ of order
$|E|+|V(R)|\ge\ell$ starting at $w$.
Take its first $\ell$ vertices and prefix whichever of $y,z$ has
color different from its last vertex.  The prefixed vertex lies
outside $K$, so the resulting path is simple and has length exactly
$\ell$.  This contradicts the choice of $\ell$ and completes the
proof.
\end{proof}

\subsection{Closing a path through the breadth-first tree}

The preceding lemma now supplies the path within one distance layer;
the tree supplies a disjoint path of the complementary length.

\begin{proof}[Proof of Theorem~\ref{clq:bound}]
Lemma~\ref{clq:triangle} gives $3\le t\le k$.
If $s=\lfloor k/2\rfloor\le3$, this is the required bound.
Suppose therefore that $s\ge4$ and that $G$ has no $K_s$.
Choose a vertex of $G$ and a breadth-first spanning tree of its
component rooted at that vertex.  Since $\alpha(G)\le a\le k$,
Proposition~\ref{clq:layer-interface} applies and gives $H$ in one
distance layer $D_i$ of this tree, where $1\le i\le s$.

In the breadth-first tree defining this layer, let $c$ be the lowest
common ancestor of all vertices of $H$.
They all have depth $i$, so their distance below $c$ is the same
integer
\[
 p=i-\operatorname{depth}(c).
\]
The graph $H$ has at least $k$ vertices, so $c$ is a proper
ancestor of its vertices, giving $1\le p\le i\le s$.
At least two child branches of $c$ meet $H$, since otherwise the
common child would be a lower common ancestor.

If $p=1$, every vertex of $H$ is adjacent to $c$.
By Lemma~\ref{clq:long-path}, the graph $H$ contains a path of
order at least
\[
 \min\{|V(H)|,2s+1\}\ge k,
\]
because $|V(H)|\ge k$ and $k\in\{2s,2s+1\}$.
Take a subpath of order exactly $k$ and join both ends to $c$.
The vertex $c$ lies above $D_i$, so this gives a simple cycle of
order $k+1$, a contradiction.

Suppose that $p\ge2$.  Partition the child branches of $c$ that
meet $H$ into two nonempty classes, and color each vertex of $H$
according to its class.  This is a nonconstant coloring.
Vertices of different colors have tree distance exactly $2p$,
and all internal vertices of their tree path have depth less than
$i$.
Set
\[
 \ell=k+1-2p.
\]
The inequalities $p\le s$, $p\ge2$, and
$k\in\{2s,2s+1\}$ give $1\le\ell\le2s-2$.
The graph $H$ has all the hypotheses of
Lemma~\ref{clq:coloured-path}: its independent-set neighborhood
bound applies in particular to independent pairs.
That lemma supplies a path in $H$ of length exactly $\ell$
with differently colored ends.
Its vertices have depth $i$, so its interior is disjoint from
the tree path between its ends.  The union of the two paths is
therefore a simple cycle of length
\[
 \ell+2p=k+1.
\]
This final contradiction proves $t\ge s$ and hence the theorem.
\end{proof}
